\begin{table}[htb]
\centering
\caption{Verificações sobre os dados: acurácia de teste da regressão logística nas nove funções $\{1, \cos x_j, \sin x_j\}$ e seus produtos, e da regra $g(x)/2 > t$, com $g(x) = \sin x_1 + \sin x_2$, no \textit{circles} (média $\pm$ desvio padrão sobre cinco sementes).}
\label{tab:verificacoes}
\begin{tabular}{lccc}
\hline
Verificação & \textit{Xor} & \textit{Moons} & \textit{Circles} \\
\hline
Controle linear (9 funções) & 0,873 $\pm$ 0,045 & 0,850 $\pm$ 0,039 & 0,963 $\pm$ 0,022 \\
$g/2 > t$, sem ruído & -- & -- & 1,000 $\pm$ 0,000 \\
$g/2 > t$, ruído 0,15 & -- & -- & 0,967 $\pm$ 0,020 \\
\hline
\end{tabular}
\\[2pt]
{\footnotesize Coordenadas já normalizadas para $[0, \pi]$, com as mesmas partições e sementes da comparação principal. A regressão logística tem hiperparâmetros fixados antes de rodar ($C = 1,0$); o limiar $t$ e o sentido da desigualdade são escolhidos no treino. Um separador nesse espaço de funções não prova que o circuito quântico realize ou aprenda os mesmos coeficientes.}
\\[2pt]
{\footnotesize Fonte: Elaborada pelo autor.}
\end{table}
